\chapter{Le produit scalaire dans le plan}\label{ch:g11:scal}

Le \omterm{def:g11:scal:dot}{produit scalaire} attache un \emph{nombre} à une paire de \omterm{def:g11:vect:vector}{vecteurs}, et
ce nombre voit la géométrie~: il s'annule exactement pour des directions
perpendiculaires, et il mesure angles et longueurs. C'est le moteur
derrière la loi des cosinus, les \omterm{def:g10:algebra:equation}{équations} de cercles, et les calculs de
perpendicularité. Le même outil, dans l'espace, est le sujet du
\cref{ch:g12:space}.

\section{Définition et premiers calculs}

\begin{definition}[Produit scalaire]\label{def:g11:scal:dot}
Le \emph{produit scalaire}\index{produit scalaire} de deux \omterm{def:g11:vect:vector}{vecteurs}
$\vec u$ et $\vec v$ est le nombre
\[
\vec u \cdot \vec v
= \tfrac12\left(\norm{\vec u + \vec v}^2 - \norm{\vec u}^2
- \norm{\vec v}^2\right),
\]
où $\norm{\vec u}$ désigne la longueur (\emph{norme}\index{norme}) de $\vec u$.
\end{definition}

\begin{theorem}[Les quatre faces du produit scalaire]\label{thm:g11:scal:faces}
Soient $\vec u\,(x, y)$ et $\vec v\,(x', y')$ des \omterm{def:g11:vect:vector}{vecteurs} non nuls, et
$\theta$ l'angle entre eux. Alors
\begin{enumerate}
  \item \emph{(\omterm{def:g10:coordgeom:system}{coordonnées})} $\vec u \cdot \vec v = xx' + yy'$~;
  \item \emph{(\omterm{def:g11:trigo:cossin}{cosinus})} $\vec u \cdot \vec v =
        \norm{\vec u}\,\norm{\vec v}\cos\theta$~;
  \item \emph{(projection)} $\vec u \cdot \vec v = \vec u \cdot \vec v'$,
        où $\vec v'$ est la projection \omterm{def:g11:scal:orthogonal}{orthogonale} de $\vec v$ sur la
        direction de $\vec u$~; si $\vec u = \vect{OA}$ et
        $\vec v = \vect{OB}$, alors
        $\vec u \cdot \vec v = \overline{OA} \times \overline{OH}$, un
        produit de longueurs algébriques, où $H$ est le pied de la
        perpendiculaire issue de $B$ à la droite $(OA)$~;
  \item \emph{(\omterm{def:g11:scal:dot}{normes})} la formule définissante de la
        \cref{def:g11:scal:dot}.
\end{enumerate}
\end{theorem}

\begin{proof}[Démonstration de 1 et 2]
\emph{1.} En \omterm{def:g10:coordgeom:system}{coordonnées}, $\norm{\vec u}^2 = x^2 + y^2$ (Pythagore) et
$\vec u + \vec v$ a pour \omterm{def:g10:coordgeom:system}{coordonnées} $(x + x', y + y')$, donc
\begin{align*}
2\,\vec u \cdot \vec v
&= (x + x')^2 + (y + y')^2 - (x^2 + y^2) - (x'^2 + y'^2)\\
&= 2xx' + 2yy',
\end{align*}
après développement des carrés et annulations.

\emph{2.} Choisir le \omterm{def:g10:coordgeom:system}{repère} pour que $\vec u$ pointe le long de l'axe des
$x$ positifs~:
\[
\vec u\,(\norm{\vec u}, 0)
\quad\text{et}\quad
\vec v\,(\norm{\vec v}\cos\theta,\ \norm{\vec v}\sin\theta)
\]
(\cref{ch:g11:trigo}). La formule~1 donne alors
$\vec u \cdot \vec v = \norm{\vec u}\,\norm{\vec v}\cos\theta + 0$. La
formule~3 lit le même calcul~: la projection de $\vec v$ sur l'axe des
$x$ a pour longueur algébrique $\norm{\vec v}\cos\theta =
\overline{OH}$.
\end{proof}

\begin{omfigure}
\begin{tikzpicture}[scale=1.25]
  \coordinate (O) at (0,0);
  \coordinate (A) at (3.2,0);
  \coordinate (B) at (55:2.4);
  \coordinate (H) at (1.3767,0);
  \draw[-Stealth, omDef, thick] (O) -- (A)
    node[below, font=\small, pos=0.85] {$\vec u$};
  \draw[-Stealth, omThm, thick] (O) -- (B)
    node[above left, font=\small, pos=0.8] {$\vec v$};
  \draw[black, dashed, thin] (B) -- (H);
  \draw[thin] (H) ++(-0.18,0) -- ++(0,0.18) -- ++(0.18,0);
  \draw[omProp, very thick] (O) -- (H);
  \fill (H) circle (1.2pt) node[below, font=\small] {$H$};
  \node[below left, font=\small] at (O) {$O$};
  \draw[->,omThm,thin] (0.55,0) arc (0:55:0.55);
  \node[omThm, font=\small] at (27:0.78) {$\theta$};
\end{tikzpicture}

{\small La vue en projection~: $\vec u \cdot \vec v = \overline{OA}
\times \overline{OH}$, où $H$ est le pied de la perpendiculaire depuis
l'extrémité de $\vec v$. Le segment orange $OH$ a pour longueur
algébrique $\norm{\vec v}\cos\theta$.}
\end{omfigure}

\begin{example}\label{ex:g11:scal:compute}
$\vec u\,(3, 1)$ et $\vec v\,(2, -4)$~:
$\vec u \cdot \vec v = 6 - 4 = 2$. Les \omterm{def:g11:scal:dot}{normes} sont
$\norm{\vec u} = \sqrt{10}$ et $\norm{\vec v} = \sqrt{20}$, donc
$\cos\theta = \frac{2}{\sqrt{10}\sqrt{20}} = \frac{2}{10\sqrt2}
= \frac{\sqrt2}{10}$~: l'angle est un peu inférieur à $82^\circ$.
\end{example}

\begin{proposition}[Règles algébriques]\label{prop:g11:scal:rules}
Pour tous \omterm{def:g11:vect:vector}{vecteurs} et tout $\lambda \in \R$~:
\[
\vec u \cdot \vec v = \vec v \cdot \vec u, \qquad
\vec u \cdot (\vec v + \vec w) = \vec u \cdot \vec v + \vec u \cdot
\vec w, \qquad
(\lambda \vec u)\cdot \vec v = \lambda\,(\vec u \cdot \vec v),
\]
\[
\vec u \cdot \vec u = \norm{\vec u}^2,
\qquad
\norm{\vec u + \vec v}^2 = \norm{\vec u}^2 + 2\,\vec u\cdot\vec v +
\norm{\vec v}^2 .
\]
\end{proposition}

\begin{proof}
Tout suit de la formule en \omterm{def:g10:coordgeom:system}{coordonnées}~: par exemple
$\vec u \cdot (\vec v + \vec w) = x(x' + x'') + y(y' + y'') =
(xx' + yy') + (xx'' + yy'')$. La dernière identité développe
$\norm{\vec u + \vec v}^2 = (\vec u + \vec v)\cdot(\vec u + \vec v)$
en utilisant les trois premières règles --- le \omterm{def:g11:scal:dot}{produit scalaire} se
comporte exactement comme un produit ordinaire, donc les identités
algébriques usuelles s'appliquent aux \omterm{def:g11:vect:vector}{vecteurs}.
\end{proof}

\section{Orthogonalité}

\begin{definition}[Vecteurs orthogonaux]\label{def:g11:scal:orthogonal}
Deux \omterm{def:g11:vect:vector}{vecteurs} sont \emph{orthogonaux}\index{vecteurs orthogonaux}, noté
$\vec u \perp \vec v$, si $\vec u \cdot \vec v = 0$.
\end{definition}

\begin{remark}
Pour des \omterm{def:g11:vect:vector}{vecteurs} non nuls, la formule~2 du \cref{thm:g11:scal:faces}
montre $\vec u \cdot \vec v = 0 \iff \cos\theta = 0 \iff$ les directions
sont perpendiculaires. Le \omterm{def:g11:vect:vector}{vecteur} nul est \omterm{def:g11:scal:orthogonal}{orthogonal} à tout.
\end{remark}

\begin{example}
$\vec u\,(3, 2)$ et $\vec v\,(-4, 6)$~: $\vec u \cdot \vec v = -12 + 12 =
0$~: \omterm{def:g11:scal:orthogonal}{orthogonaux}. En général, $\vec u\,(a, b)$ est toujours \omterm{def:g11:scal:orthogonal}{orthogonal} à
$\vec v\,(-b, a)$ --- une rotation d'un quart de tour.
\end{example}

\section{La loi des cosinus}

\begin{theorem}[Loi des cosinus]\label{thm:g11:scal:cosines}
\index{loi des cosinus}
Dans un triangle $ABC$ de côtés $a = BC$, $b = CA$, $c = AB$ et d'angle
$\widehat A$ au sommet $A$~:
\[
a^2 = b^2 + c^2 - 2bc\,\cos\widehat A .
\]
\end{theorem}

\begin{proof}
Écrire $\vect{BC} = \vect{BA} + \vect{AC} = \vect{AC} - \vect{AB}$ et
\omterm{def:g10:algebra:expand}{développer} avec la \cref{prop:g11:scal:rules}~:
\[
a^2 = \norm{\vect{BC}}^2
= \norm{\vect{AC}}^2 - 2\,\vect{AC}\cdot\vect{AB} + \norm{\vect{AB}}^2
= b^2 + c^2 - 2\,\vect{AB}\cdot\vect{AC}.
\]
Par la formule du \omterm{def:g11:trigo:cossin}{cosinus},
$\vect{AB}\cdot\vect{AC} = c\,b\,\cos\widehat A$.
\end{proof}

\begin{remark}
Lorsque $\widehat A$ est un angle droit, $\cos\widehat A = 0$ et la loi
des cosinus se réduit au théorème de Pythagore $a^2 = b^2 + c^2$~:
c'est Pythagore avec un terme correctif pour les angles non droits.
\end{remark}

\begin{omfigure}
\begin{tikzpicture}[scale=1.15]
  \coordinate (A) at (0,0);
  \coordinate (B) at (4,0);
  \coordinate (C) at (1.2,2.2);
  \draw[omDef, thick] (A) -- (B) -- (C) -- cycle;
  \node[below left, font=\small] at (A) {$A$};
  \node[below right, font=\small] at (B) {$B$};
  \node[above, font=\small] at (C) {$C$};
  \node[below, font=\small] at (2,0) {$c$};
  \node[above left, font=\small] at (0.6,1.1) {$b$};
  \node[above right, font=\small] at (2.6,1.1) {$a$};
  \draw[->,omThm,thin] (0.7,0) arc (0:61.4:0.7);
  \node[omThm, font=\small] at (30:1.0) {$\widehat A$};
\end{tikzpicture}

{\small La loi des cosinus~: connaître deux côtés $b$, $c$ et l'angle
entre eux détermine le troisième côté $a$.}
\end{omfigure}

\begin{example}\label{ex:g11:scal:cosines}
Deux côtés d'un triangle mesurent $b = 5$ et $c = 8$, enfermant un angle
$\widehat A = \frac\pi3$. Alors
\[
a^2 = 25 + 64 - 2 \times 5 \times 8 \times \frac12 = 89 - 40 = 49,
\]
donc $a = 7$.
\end{example}

\section{Applications~: cercles et droites perpendiculaires}

\begin{proposition}[Cercle de diamètre donné]\label{prop:g11:scal:diameter}
Un point $M$ est sur le cercle de diamètre $[AB]$ si et seulement si
\[
\vect{MA} \cdot \vect{MB} = 0 .
\]
\end{proposition}

\begin{proof}
Si $M \neq A, B$, la condition dit que l'angle en $M$ dans le triangle
$AMB$ est droit, et les points qui voient un segment sous un angle droit
sont exactement ceux du cercle de ce diamètre. Voici une preuve directe~:
soit $O$ le \omterm{prop:g10:coordgeom:midpoint}{milieu} de $[AB]$ et $r = \frac{AB}{2}$. Alors
$\vect{MA} = \vect{MO} + \vect{OA}$ et
$\vect{MB} = \vect{MO} + \vect{OB} = \vect{MO} - \vect{OA}$, donc
\[
\vect{MA}\cdot\vect{MB}
= \norm{\vect{MO}}^2 - \norm{\vect{OA}}^2 = MO^2 - r^2,
\]
qui s'annule exactement lorsque $MO = r$, c'est-à-dire lorsque $M$ est
sur le cercle de centre $O$ et de rayon $r$.
\end{proof}

\begin{proposition}[Vecteur normal d'une droite]\label{prop:g11:scal:normal}
Le \omterm{def:g11:vect:vector}{vecteur} $\vec n\,(a, b)$ est \omterm{def:g11:scal:orthogonal}{orthogonal} à tout \omterm{def:g11:vect:direction}{vecteur directeur} de la
droite $d : ax + by + c = 0$~; on dit que $\vec n$ est un \emph{vecteur
normal}\index{vecteur normal} de $d$. Deux droites sont perpendiculaires
exactement lorsque leurs \omterm{def:g11:vect:vector}{vecteurs} normaux (ou directeurs) sont
\omterm{def:g11:scal:orthogonal}{orthogonaux}.
\end{proposition}

\begin{proof}
Un \omterm{def:g11:vect:direction}{vecteur directeur} de $d$ est $\vec u\,(-b, a)$
(\cref{thm:g11:vect:cartesian}), et
$\vec n \cdot \vec u = a(-b) + b(a) = 0$.
\end{proof}

\begin{method}[Utiliser les vecteurs normaux]\label{met:g11:scal:normal}
\begin{itemize}
  \item Droite par $A$ de normale $\vec n\,(a,b)$~: écrire
        $\vect{AM} \cdot \vec n = 0$, qui se développe en
        $ax + by + c = 0$ avec $c$ déterminé par $A$.
  \item Perpendicularité de $d: ax+by+c = 0$ et $d': a'x+b'y+c' = 0$~:
        vérifier $aa' + bb' = 0$.
\end{itemize}
\end{method}

\begin{example}\label{ex:g11:scal:perpline}
La droite par $P(2, -1)$ perpendiculaire à $d : 3x - y + 4 = 0$~: un
\omterm{def:g11:vect:direction}{vecteur directeur} de $d$ est $(1, 3)$, qui sert de \omterm{def:g11:vect:vector}{vecteur}
\emph{normal} de la perpendiculaire. \omterm{def:g10:algebra:equation}{Équation}~: $x + 3y + c = 0$ avec
$2 - 3 + c = 0$, donc $c = 1$~: la droite est $x + 3y + 1 = 0$.
\end{example}

\begin{example}[Équation d'un cercle]\label{ex:g11:scal:circle}
Le cercle de centre $\Omega(1, 2)$ et de rayon $3$ est l'ensemble des
points $M(x,y)$ tels que $\Omega M^2 = 9$~:
\[
(x - 1)^2 + (y - 2)^2 = 9 .
\]
Réciproquement, $x^2 + y^2 - 2x - 4y - 4 = 0$ se réécrit, en complétant
les carrés (\cref{ch:g11:quad}), comme $(x-1)^2 + (y-2)^2 = 9$~: le même
cercle.
\end{example}

\section{Exercices}

\begin{exercise}[$\star$]\label{exo:g11:scal:1}
Calculer $\vec u \cdot \vec v$ pour
\[
\vec u\,(2, 5),\ \vec v\,(3, -1); \qquad
\vec u\,(1, -4),\ \vec v\,(8, 2); \qquad
\norm{\vec u} = 3,\ \norm{\vec v} = 4,\ \theta = \frac{\pi}{3} .
\]
\end{exercise}

\begin{exercise}[$\star$]\label{exo:g11:scal:2}
$ABC$ est un triangle équilatéral de côté $6$. Calculer
$\vect{AB} \cdot \vect{AC}$ et $\vect{AB} \cdot \vect{BC}$.
\end{exercise}

\begin{exercise}[$\star$]\label{exo:g11:scal:3}
Pour quelle valeur de $t$ les \omterm{def:g11:vect:vector}{vecteurs} $\vec u\,(3, t)$ et
$\vec v\,(t - 8, 1)$ sont-ils \omterm{def:g11:scal:orthogonal}{orthogonaux}~? Et pour quelle(s) valeur(s)
de $t$ le \omterm{def:g11:vect:vector}{vecteur} $\vec u\,(t, 2)$ est-il \omterm{def:g11:scal:orthogonal}{orthogonal} à lui-même moins
$\vec v\,(4, 2)$, c'est-à-dire $\vec u \perp (\vec u - \vec v)$~?
\end{exercise}

\begin{exercise}[$\star$]\label{exo:g11:scal:4}
Calculer l'angle $\theta$ entre $\vec u\,(1, 2)$ et $\vec v\,(3, 1)$
(donner $\cos\theta$ exactement, puis $\theta$ au degré le plus proche).
\end{exercise}

\begin{exercise}[$\star\star$]\label{exo:g11:scal:5}
Un triangle a pour côtés $b = 4$, $c = 6$ et pour angle
$\widehat A = \frac{2\pi}{3}$ entre eux. Calculer le troisième côté $a$.
Puis, dans un triangle de côtés $a = 7$, $b = 5$, $c = 4$, calculer
$\cos\widehat A$ et décider si $\widehat A$ est aigu ou obtus.
\end{exercise}

\begin{exercise}[$\star\star$]\label{exo:g11:scal:6}
Soient $A(1, 1)$, $B(5, -1)$ et $C(3, 5)$. Calculer
$\vect{AB} \cdot \vect{AC}$~; l'angle $\widehat A$ est-il droit~?
Calculer les trois longueurs de côtés et vérifier avec la loi des
cosinus.
\end{exercise}

\begin{exercise}[$\star\star$]\label{exo:g11:scal:7}
Donner une \omterm{def:g10:algebra:equation}{équation} du cercle de centre $\Omega(-2, 3)$ et de rayon $5$.
Puis déterminer le centre et le rayon du cercle
\[
x^2 + y^2 - 6x + 2y - 6 = 0 .
\]
\end{exercise}

\begin{exercise}[$\star\star$]\label{exo:g11:scal:8}
Soient $A(-3, 0)$ et $B(0, 2)$. En utilisant la
\cref{prop:g11:scal:diameter}, trouver une \omterm{def:g10:algebra:equation}{équation} du cercle de
diamètre $[AB]$, et vérifier que l'origine $O(0,0)$ y est.
\end{exercise}

\begin{exercise}[$\star\star$]\label{exo:g11:scal:9}
Trouver une \omterm{def:g10:algebra:equation}{équation} de la droite passant par $P(1, 3)$ perpendiculaire
à $d : 2x + y - 7 = 0$, et les \omterm{def:g10:coordgeom:system}{coordonnées} du pied $H$ de la
perpendiculaire (l'\omterm{def:g10:numbers:interunion}{intersection} des deux droites). En déduire la
distance de $P$ à $d$.
\end{exercise}

\begin{exercise}[$\star\star$]\label{exo:g11:scal:10}
En utilisant $\norm{\vec u + \vec v}^2 = \norm{\vec u}^2 +
2\vec u\cdot\vec v + \norm{\vec v}^2$ et son analogue pour
$\vec u - \vec v$, prouver l'\emph{identité du parallélogramme}~:
\[
\norm{\vec u + \vec v}^2 + \norm{\vec u - \vec v}^2
= 2\norm{\vec u}^2 + 2\norm{\vec v}^2 ,
\]
et l'interpréter dans un parallélogramme (diagonales contre côtés).
\end{exercise}

\begin{exercise}[$\star\star\star$]\label{exo:g11:scal:11}
Soient $A$ et $B$ deux points avec $AB = 4$, et $I$ le \omterm{prop:g10:coordgeom:midpoint}{milieu} de
$[AB]$.
\begin{enumerate}
  \item Montrer que pour tout point $M$~:
        $\vect{MA}\cdot\vect{MB} = MI^2 - 4$.
        (Insérer $I$~: $\vect{MA} = \vect{MI} + \vect{IA}$.)
  \item En déduire l'ensemble des points $M$ tels que
        $\vect{MA}\cdot\vect{MB} = 12$.
\end{enumerate}
\end{exercise}
